\documentclass{article}
\usepackage{amsmath, amssymb}

\begin{document}

\section*{Pushing--Medium Gravity: Hamiltonian Formulation v1}

This document presents the Hamiltonian formulation of the Pushing--Medium (PM)
medium derived from the PM-fluid Lagrangian. The Hamiltonian density describes
the total energy of the PM medium in canonical variables and provides the basis
for stability analysis and canonical evolution.

\section*{1. PM Lagrangian Density}

The PM-fluid Lagrangian on a flat background $\eta_{\mu\nu} =
\mathrm{diag}(-1,1,1,1)$ is


\[
\mathcal{L}
=
\frac{1}{2c_\phi^2}(\partial_t \phi)^2
- \frac{1}{2}|\nabla \phi|^2
+
\frac{1}{2c_u^2}|\partial_t \mathbf u|^2
- \frac{1}{2}|\nabla \mathbf u|^2
+
\phi\,\nabla\cdot(\phi\mathbf u)
+
\phi\,S_\phi
+
\mathbf u\cdot\mathbf S_u.
\]



For the Hamiltonian we omit the explicit source terms
$\phi S_\phi$ and $\mathbf u\cdot\mathbf S_u$.

\section*{2. Canonical Momenta}

\subsection*{2.1 Scalar Field}



\[
\pi_\phi
=
\frac{\partial \mathcal{L}}{\partial(\partial_t \phi)}
=
\frac{1}{c_\phi^2}\,\partial_t \phi.
\]



\subsection*{2.2 Vector Field}



\[
\pi_{u_i}
=
\frac{\partial \mathcal{L}}{\partial(\partial_t u_i)}
=
\frac{1}{c_u^2}\,\partial_t u_i.
\]



These define the canonical phase space


\[
(\phi, \pi_\phi), \qquad (u_i, \pi_{u_i}).
\]



\section*{3. Hamiltonian Density}

The Hamiltonian density is


\[
\mathcal{H}
=
\pi_\phi\,\partial_t \phi
+
\pi_{u_i}\,\partial_t u_i
- \mathcal{L}.
\]



Substituting the momenta and eliminating time derivatives:



\[
\partial_t \phi = c_\phi^2 \pi_\phi,
\qquad
\partial_t u_i = c_u^2 \pi_{u_i}.
\]



Thus:


\[
\mathcal{H}
=
\frac{1}{2}c_\phi^2 \pi_\phi^2
+ \frac{1}{2}|\nabla \phi|^2
+
\frac{1}{2}c_u^2 |\pi_{u}|^2
+ \frac{1}{2}|\nabla \mathbf u|^2
- \phi\,\nabla\cdot(\phi\mathbf u).
\]



This is the Hamiltonian density of the PM medium.

\section*{4. Interpretation of Terms}

\begin{itemize}
    \item $c_\phi^2 \pi_\phi^2 / 2$ is the kinetic energy of scalar PM waves.
    \item $|\nabla \phi|^2 / 2$ is the gradient (tension) energy of the scalar field.
    \item $c_u^2 |\pi_u|^2 / 2$ is the kinetic energy of vector PM waves.
    \item $|\nabla \mathbf u|^2 / 2$ is the gradient energy of the flow field.
    \item $-\phi\,\nabla\cdot(\phi\mathbf u)$ is the compressive interaction energy between $\phi$ and $\mathbf u$.
\end{itemize}

\section*{5. Hamilton’s Equations}

The canonical evolution equations are:



\[
\partial_t \phi = \frac{\delta \mathcal{H}}{\delta \pi_\phi}
= c_\phi^2 \pi_\phi,
\]





\[
\partial_t \pi_\phi = -\frac{\delta \mathcal{H}}{\delta \phi}
= \nabla^2 \phi + \nabla\cdot(\phi\mathbf u) + \phi\,\nabla\cdot\mathbf u,
\]





\[
\partial_t u_i = \frac{\delta \mathcal{H}}{\delta \pi_{u_i}}
= c_u^2 \pi_{u_i},
\]





\[
\partial_t \pi_{u_i} = -\frac{\delta \mathcal{H}}{\delta u_i}
= \nabla^2 u_i + \partial_i(\phi^2).
\]



These reproduce the PM-fluid equations in the absence of external sources.

\section*{6. Stability and Positivity}

The Hamiltonian density is positive definite in the kinetic and gradient terms:


\[
\frac{1}{2}c_\phi^2 \pi_\phi^2
+ \frac{1}{2}|\nabla \phi|^2
+
\frac{1}{2}c_u^2 |\pi_u|^2
+ \frac{1}{2}|\nabla \mathbf u|^2 \ge 0.
\]



The interaction term


\[
-\phi\,\nabla\cdot(\phi\mathbf u)
\]


is linear in $\mathbf u$ and quadratic in $\phi$ and does not destabilize the Hamiltonian for small flows. This confirms that the PM medium is stable and bounded below in the linear regime validated by numerical tests.

\section*{7. Deformation Energy and Matter Stability}

The homogeneous (spatially uniform, static) sector of the Hamiltonian reduces
to the \emph{deformation energy} density, which governs the stability of
matter-like configurations in the PM medium.

\subsection*{7.1 Landau--Ginzburg Form}

For a uniform, static medium ($\mathbf u = 0$, $\pi_\phi = 0$,
$\nabla\phi = 0$) the Hamiltonian density reduces to the potential energy
stored in the scalar field.  Retaining the leading non-linear self-interaction
of $\phi$ generated by the coupling $\phi\,\nabla\cdot(\phi\mathbf u)$
(adiabatically integrating out the flow field driven by $\nabla\phi$), the
effective deformation energy density takes the Landau--Ginzburg form:


\[
U(\phi)
=
\varepsilon_0\!\left(\phi^2 - \frac{\phi^3}{3}\right),
\qquad
\varepsilon_0 = \rho_{\mathrm{nuc}}\,c^2.
\]


The first and second field-derivatives are:


\[
U'(\phi) = \varepsilon_0(2\phi - \phi^2),
\qquad
U''(\phi) = 2\varepsilon_0(1 - \phi).
\]


\subsection*{7.2 Stability Regions}

A matter configuration is stable when the deformation energy is locally convex:


\[
U''(\phi) > 0
\quad\Longleftrightarrow\quad
\phi < 1.
\]


The stability regimes are:

\begin{itemize}
    \item $\phi < 1$: $U'' > 0$ -- stable matter configuration.
    \item $\phi = 1$: $U'' = 0$ -- critical point (onset of instability).
    \item $\phi > 1$: $U'' < 0$ -- unstable; matter dissolves into excitations.
\end{itemize}

\subsection*{7.3 Critical Compression}

The critical compression $\phi_{\mathrm{crit}}$ is the smallest positive root of $U'' = 0$:


\[
\phi_{\mathrm{crit}} = 1.
\]


The corresponding PM refractive index, density, and pressure at the transition are:


\[
n_{\mathrm{crit}} = e^{\phi_{\mathrm{crit}}} = e,
\]


\[
\rho_{\mathrm{crit}}
= \rho_{\mathrm{nuc}}\,e^{\phi_{\mathrm{crit}}}
= e\,\rho_{\mathrm{nuc}}
\approx 6.25\times10^{17}\;\mathrm{kg\,m^{-3}},
\]


\[
P_{\mathrm{crit}}
= \frac{\rho_{\mathrm{crit}}\,c^2}{3}
\approx 1.87\times10^{34}\;\mathrm{Pa}.
\]


This corresponds to approximately $2.7\,\rho_{\mathrm{nuc}}$, within the broad
bandwidth associated with the onset of nuclear-matter failure.

\subsection*{7.4 Physical Interpretation}

The cubic term $-\phi^3/3$ in $U(\phi)$ is the leading non-linear correction
arising from the PM coupling $\phi\,\nabla\cdot(\phi\mathbf u)$ in the
Lagrangian.  Its sign ensures the energy landscape tilts downward for
$\phi > 1$, making configurations at high compression energetically
unfavourable and mechanically unstable.  This is the PM analogue of the
spinodal instability in classical thermodynamics.

Above the critical compression the medium transitions into a radiation-like,
ultra-relativistic state ($P = \rho c^2/3$) with no net matter-like order
parameter.  This is the PM matter--energy transition described in
\texttt{critical\_state\_computation.md}.

\subsection*{7.5 Implementation}

The helpers \texttt{pm\_deformation\_energy}, \texttt{pm\_deformation\_energy\_deriv1},
\texttt{pm\_deformation\_energy\_deriv2}, \texttt{pm\_density\_from\_phi}, and
\texttt{pm\_pressure\_from\_phi} in
\texttt{src/pushing\_medium/critical\_state.py} implement these relations.
The function \texttt{compute\_critical\_state()} finds $\phi_{\mathrm{crit}}$
by scanning $U''(\phi)$ and refining with 64-step bisection, returning a
\texttt{CriticalState} dataclass.  Full numerical verification is in
\texttt{tests/test\_critical\_state.py} (11 tests, all passing).

\section*{8. Summary}

The Hamiltonian formulation provides:

\begin{itemize}
    \item canonical variables $(\phi, \pi_\phi)$ and $(u_i, \pi_{u_i})$,
    \item a positive-definite Hamiltonian density,
    \item canonical evolution equations equivalent to the PM-fluid equations,
    \item a clear decomposition of PM energy into scalar, vector, gradient, and interaction components,
    \item a deformation-energy landscape $U(\phi)$ whose curvature criterion identifies the PM matter/energy transition at $\phi_{\mathrm{crit}} = 1$.
\end{itemize}

This structure completes the foundational mechanics of the PM medium and
prepares the theory for further analysis, including wave energetics, conserved
charges, potential quantization, and the microphysics of compact objects.

\end{document}

